\begin{table}[ht]
\centering
\caption{Backcast residential population of downstream community water systems: summary statistics.}
\label{tab:backcast_popsum}
\begin{adjustbox}{max width=\textwidth}
\begin{tabular}{lrrrrrrrrr}
\toprule
 & N & Mean & SD & p10 & p25 & Median & p75 & p90 & Max \\
\midrule
\multicolumn{10}{l}{\textit{Panel A: Pooled system $\times$ year observations, 1990--2024}} \\
All systems & 12,810 & 53,947 & 138,869 & 24 & 166 & 3,993 & 56,755 & 128,487 & 1,340,627 \\
\quad Service-area tier & 8,365 & 10,234 & 42,706 & 11 & 61 & 480 & 3,565 & 20,578 & 537,146 \\
\quad County tier & 4,445 & 136,209 & 204,417 & 29,788 & 42,892 & 78,797 & 144,614 & 319,842 & 1,340,627 \\
\midrule
\multicolumn{10}{l}{\textit{Panel B: Cross-section of all recovered systems, selected years}} \\
1990 & 366 & 52,800 & 141,423 & 16 & 127 & 3,815 & 53,749 & 124,524 & 1,336,449 \\
1995 & 366 & 53,809 & 141,668 & 17 & 129 & 3,947 & 53,681 & 125,666 & 1,322,172 \\
2000 & 366 & 53,497 & 138,866 & 19 & 150 & 4,060 & 54,487 & 123,762 & 1,281,666 \\
2005 & 366 & 53,575 & 136,666 & 18 & 151 & 4,019 & 54,330 & 122,290 & 1,236,731 \\
\bottomrule
\end{tabular}
\end{adjustbox}
\begin{minipage}{\textwidth}\footnotesize\raggedright
\vspace{0.5em}
\textit{Notes:} The unit is the residential population inside a system's exposure polygon ($\hat{G}_{i,t}$), not the population the utility reports serving. The distribution is right-skewed and bimodal because the two geography tiers differ in scale: service-area polygons delineate the utility's own footprint, whereas county-tier systems inherit their entire county's population and so overstate the population attributable to any one system. Panel B holds the set of systems fixed at all recovered systems, isolating population change from roster entry and exit.
\end{minipage}
\end{table}
